\newpage
\pagestyle{fancy}

\section{Criteri di progetto degli impianti}
Tutti i materiali e le apparecchiature utilizzati devono essere di alta qualità, prodotti da aziende affidabili, ben lavorati e adatti all'uso previsto, resistendo a sollecitazioni meccaniche, corrosione, calore, umidità e acque meteoriche (per installazione all’esterno). Devono garantire lunga durata, facilità di ispezione e manutenzione. 
È obbligatorio l'uso di componenti con marcatura CE e, se disponibile, marchio IMQ o equivalente europeo. I componenti senza marcatura CE devono avere una dichiarazione di conformità del costruttore ai requisiti di sicurezza delle normative CEI, UNI o IEC.

\subsection{Dimensionamento delle linee}
Le linee elettriche sono calcolate mediante l’utilizzo dei seguenti criteri progettuali:
\begin{itemize}
    \item La corrente di impiego (Ib) è calcolata considerando la potenza nominale delle apparecchiature elettriche. La tensione di alimentazione è pari a 230V per le utenze monofase, 400 V per le utenze trifase. Fattore di potenza pari a 0,99 per le linee di alimentazione delle prese di ricarica.
    \item La corrente nominale della protezione (In), definita dal costruttore, è considerata come la corrente che un interruttore può sopportare per un tempo indefinito senza che quest’ultimo subisca alcun danno;
    \item La portata del cavo (Iz) è calcolata utilizzando le tabelle CEI UNEL 35024 e 35026, tenendo conto delle condizioni di posa, il tipo di isolante del cavo e della temperatura ambiente.
\end{itemize}
I cavi di alimentazione sono stati inoltre dimensionati in modo da non subire danneggiamento causato da sovraccarichi e cortocircuiti mediante il coordinamento con la corrente nominale (In) del dispositivo di protezione a monte in linea a quanto descritto nei successivi paragrafi (Sovraccarichi e cortocircuiti di cui ai paragrafi 3.5.1 e 3.5.2)


\subsection{Calcolo della sezione del cavo in funzione della corrente di impiego (Ib)}
Nota la potenza assorbita dall’utenza, la corrente d’impiego (Ib), ossia la corrente assorbita da un carico, può essere calcolata come: 
\[ I_b = \frac{K_u \cdot P}{k \cdot V_n \cdot \cos\varphi} \]

dove:
\begin{itemize}
\item $k$ è pari a 1 per i circuiti monofase, $\sqrt{3}$ per i circuiti trifase;
\item $K_u$ è il coefficiente di utilizzazione della potenza nominale del carico e indica la percentuale di utilizzo dell'utenza;
\item $P$ è la potenza totale dell'utenza, calcolata in Watt [W];
\item $V_n$ è la tensione nominale del sistema [V].
\end{itemize}

Determinata la corrente di impiego per ogni utenza, è possibile dimensionare il cavo per alimentarla con una corrente di portata $I_z$ superiore ad $I_b$.

\subsection{Caduta di tensione}
Oltre ai fattori precedentemente indicati, dopo aver determinato la sezione del cavo in funzione della corrente d’impiego, è necessario verificare la caduta di tensione lungo la linea con la formula che segue:
La caduta di tensione è calcolata con la seguente formula:
\[
\Delta V = K \cdot (R \cos\varphi + X \sin\varphi) \cdot L \cdot I
\]
dove:
\begin{itemize}
\item $K$ è un coefficiente pari a 2 per le linee monofase (230\,V) e a $\sqrt{3}$ per le linee trifase (400\,V);
\item $R$ e $X$ sono rispettivamente la resistenza e la reattanza unitaria dei cavi, espresse in [$\Omega$/km];
\item $\cos\varphi$ e $\sin\varphi$ sono legati al fattore di potenza;
\item $L$ è la lunghezza della linea in metri [m];
\item $I$ è la corrente di impiego dell'apparecchiatura alimentata dal conduttore su cui si verifica la caduta di tensione.
\end{itemize}

La caduta di tensione percentuale viene calcolata con la formula:
\[
\Delta V_{\%} = \frac{\Delta V}{V} \cdot 100
\]
dove $\Delta V$ è la caduta di tensione ricavata dalla formula precedente e $V$ è la tensione nominale del sistema (230\,V per utenze monofase, 400\,V per utenze trifase).

La caduta di tensione percentuale, dal punto di fornitura (inizio linea) fino all'utenza finale (fine linea), non deve essere superiore al 4\%, come raccomandato dalla Norma CEI 64-8 art. 525.

\subsection{Sezione e tipologia dei cavi utilizzati}
I cavi utilizzati nella realizzazione dell’impianto dell’intervento in oggetto sono conformi ai requisiti previsti dal regolamento dell’Unione Europea 305/2011 (prodotti da costruzione CPR), all’unificazione UNEL e alle norme costruttive stabilite dal Comitato Elettrotecnico Italiano (CEI). \\
Determinata la sezione dei conduttori di fase come descritto nei capitoli precedenti, è necessario considerare il dimensionamento dei conduttori di neutro e del conduttore di protezione (PE). \\
La norma CEI 64-8/5 par 543.1.2 tabella 54F stabilisce la relazione tra le sezioni dei conduttori di protezione e dei conduttori di fase:

\begin{itemize}
\item La sezione del conduttore di neutro non dovrà essere inferiore a quella dei corrispondenti conduttori di fase. Per i circuiti polifase con sezione superiore a 16\,mm$^2$, la sezione del neutro è calcolata come la metà di quella dei conduttori di fase, con un minimo di 16\,mm$^2$;
\item La sezione del conduttore di protezione (PE) è determinata con riferimento alla tabella 54F della norma CEI 64-8, in particolare:

\begin{itemize}
\item Per una sezione del conduttore di fase $S_f \leq 16$\,mm$^2$, allora $S_{PE} = S_f$;
\item Per una sezione del conduttore di fase $16$\,mm$^2 < S_f \leq 35$\,mm$^2$, allora $S_{PE} = 16$\,mm$^2$;
\item er una sezione del conduttore di fase $S_f > 35$\,mm$^2$, allora $S_{PE} = S_f/2$.
\end{itemize}
\end{itemize}

Qualora il conduttore di protezione non faccia parte del cavo di alimentazione, le sue dimensioni sono determinate come sopra riportato, con valore minimo di 2,5\,mm$^2$ in rame o 16\,mm$^2$ in alluminio se è prevista protezione meccanica, e di 4\,mm$^2$ in rame o 16\,mm$^2$ in alluminio in assenza di protezione meccanica.


\subsubsection{Tipologia dei cavi}
I cavi impiegati nel seguente progetto sono del tipo:
\begin{itemize}
    \item In linea con le normative antincendio citate, cavi conformi ai requisiti del Regolamento UE 305/2011 (cavi CPR) aventi classe di reazione al fuoco \texttt{Cca-s1b,d1,a1}, tensione nominale $U_0$/U 0,6/1\,kV, con conduttore flessibile in rame rosso ricotto, a doppio isolamento costituito da isolante elastomerico reticolato (HEPR) di qualità G16, e guaina termoplastica LSZH di qualità M16, unipolari/multipolari, sigla di designazione \textbf{\{tipoDiCavo\}}, conformi alla norma CEI UNEL 35324.
    \item Cavo di categoria 0 (ad esempio cavi LAN), utilizzati per alimentazione sistemi a bassissima tensione (non superiore a 50\,V in corrente alternata e 120\,V in corrente continua). Inoltre, in linea a quanto specificato nella norma CEI-UNEL 36762, i cavi devono riportare la marcatura \textbf{``CEI-UNEL 36762 C-4 ($U_0 = 400$\,V)''} per indicare la conformità alle condizioni di coesistenza con cavi di potenza aventi tensione nominale verso terra fino a 400\,V in corrente alternata.
\end{itemize}


\subsubsection{Posa dei cavi}
Negli impianti in oggetto sono previste le seguenti tipologie di posa dei cavi e dei conduttori in generale:
\begin{itemize}
    \item Entro tubazioni interrate per le distribuzioni esterne: si dovranno prevedere pozzetti di ispezione in muratura a una distanza massima di 20 m l’uno dall’altro in modo da consentire l’infilaggio e lo sfilaggio dei cavi;
    \item Su passerelle portacavi orizzontali, verticali, inclinate: i cavi dovranno essere fissati alle passerelle mediante legature che mantengano fissati i cavi alla struttura, in particolare nei tratti verticali e inclinati dove le legature dovranno essere adatte a sostenere il peso dei cavi stessi;
    \item Entro tubazioni a vista o incassate: in questa tipologia di posa le dimensioni interne dei tubi devono essere tali da assicurare un comodo infilaggio e sfilaggio dei cavi contenuti e la superficie interna del tubo dovrà essere sufficientemente liscia perché l’infilaggio dei cavi non danneggi la guaina e/o l’isolante di questi.
\end{itemize}
Qualora fosse necessario, \textbf{sarà utile ripristinare la compartimentazione degli attraversamenti nei muri e solette} mediante l’utilizzo di barriere resistenti al fuoco tale da conservare il grado REI della struttura.
La norma CEI-UNEL 36762 fornisce indicazioni sulla posa dei cavi di categoria 0, in particolare sull’installazione in un’unica conduttura con cavi di categoria I (bassa tensione, non superiore a 1000 V in corrente alternata e 1500 V in corrente continua). Secondo tale norma i cavi di categoria 0 possono coesistere nella stessa conduttura con cavi di categoria I, senza l’utilizzo di setto separatore, nelle condizioni seguenti: 
\begin{itemize}
    \item \textbf{Ogni cavo sia isolato per la tensione più elevata presente}: l’isolamento dei cavi LAN deve essere idoneo a coesistere con la tensione nominale dei sistemi di potenza;
    \item \textbf{Ogni anima di un cavo multipolare sia isolata per la tensione più elevata presente nel cavo.}
\end{itemize}
Qualora le condizioni seguenti non siano soddisfatte è necessario utilizzare due condutture separate o inserire un setto separatore in modo tale che i cavi di categoria I siano separati da quelli di categoria 0.


\subsubsection{Colorazione dei conduttori}
I conduttori saranno identificati con i colori definiti dalle tabelle di unificazione CEI-UNEL 00722-74 e 00712; in particolare:
\begin{itemize}
    \item I conduttori di protezione saranno contraddistinti esclusivamente con il colore giallo/verde;
    \item I conduttori di neutro saranno contraddistinti con il colore blu;
    \item I conduttori di fase saranno contraddisti con i colori marrone, grigio e nero.
\end{itemize}



\subsection{Protezioni dalle sovracorrenti}
La protezione delle linee di alimentazione dalle sovracorrenti (sovraccarichi e cortocircuiti) viene assicurata da interruttori automatici di tipo magnetotermico dimensionati in modo da garantire che le curve di intervento I – t (tempo – corrente) si mantengano al di sotto delle curve caratteristiche dei cavi da essi protetti.
Per garantire la loro funzione gli interruttori automatici dovranno:
\begin{itemize}
    \item Interrompere sia la corrente di sovraccarico che quella di cortocircuito prima che esse possano provocare un aumento della temperatura nel conduttore tale da danneggiarne l’isolamento;
    \item Essere installati all’origine di ogni circuito e di ogni derivazione avente portate differenti (diverse sezioni dei conduttori, diversi tipi di posa, diverso tipo di isolamento);
    \item Avere un PdI (potere di interruzione) maggiore della massima corrente di cortocircuito presunta che si può avere nel punto di installazione del dispositivo.
\end{itemize}

\subsubsection{Sovraccarichi}
Per garantire la protezione dai sovraccarichi deve essere rispettata la relazione seguente, indicata dalla Norma CEI 64-8 art. 433.2:
\[ I_b \leq I_n \leq I_z \]
\[ I_f \leq 1{,}45 \cdot I_z \]

dove: \\
- Ib è la corrente di impiego del circuito stimata in base al carico elettrico dell’utenza. \\
- In è la corrente nominale del dispositivo di protezione. \\
- Iz è la portata del cavo ricavata tenendo conto delle condizioni di posa dello stesso cavo. \\
- If è la corrente che assicura il funzionamento del dispositivo di protezione entro il tempo convenzionale. \\
La seconda delle due disuguaglianze citate prima è soddisfatta in automatico nel caso in cui vengano utilizzati interruttori automatici conformi alle norme CEI 23-3 e CEI 17-5. \\
Nel dimensionamento dell’impianto oggetto del seguente intervento sono stati considerati gli interruttori avente corrente nominale (In) di valore tale da soddisfare la prima disuguaglianza.



\subsubsection{Cortocircuiti}
La protezione dei cortocircuiti viene espressamente delineata nell’articolo 434.3 della Norma CEI 64-8 che stabilisce che il dispositivo di protezione deve interrompere le correnti di cortocircuito in modo tale da garantire una completa protezione del conduttore, cosicché da non permettere il raggiungimento di una temperatura pericolosa, secondo la seguente relazione: 
\[ I^2 \cdot t \leq K^2 \cdot S^2 \]

dove: \\
- $I^2 t$ è l'integrale di Joule per la durata del cortocircuito. \\
- I è il valore efficace, misurata in Ampere, della corrente di guasto che percorre il dispositivo di interruzione per guasto di impedenza trascurabile. \\
- t è il tempo di intervento, in secondi, del dispositivo di protezione. \\
- K è il coefficiente tipico del cavo che dipende dall’isolante e dal materiale conduttore che costituiscono il cavo e dalla temperatura. \\
- S è la sezione del cavo in $mm^2$.

La corrente massima di cortocircuito che di solito si ha nei guasti trifase, si calcola nelle condizioni che originano i valori più elevati:
\begin{itemize}
    \item All’inizio della linea, quando l’impedenza a monte è minima;
    \item Considerando il guasto di tutti i conduttori quando la linea è costituita da più cavi in parallelo.
\end{itemize}

La corrente minima di cortocircuito che di solito si ha nei guasti monofase o bifase, si calcola nelle condizioni che originano i valori più bassi:
\begin{itemize}
    \item In fondo alla linea quando l’impedenza a monte è massima;
    \item Considerando guasti che riguardano un solo conduttore per più cavi in parallelo.
\end{itemize}


\subsection{Protezione dai contatti indiretti}
La protezione contro i contatti indiretti deve essere realizzata mediante l’utilizzo dell’interruzione automatica dell’alimentazione e/o mediante l’uso di componenti a doppio isolamento. Per la prima il tipo di protezione cambia a seconda del sistema di alimentazione, TT o TN. 

\subsubsection{Interruzione automatica dell’alimentazione in un sistema TT}
Nei sistemi TT per garantire la protezione dai contatti indiretti si utilizza un dispositivo di protezione di tipo differenziale che viene scelto mediante la formula seguente (Norma CEI 64-8 art. 413.1.4.2):
\[ I_{dn} \leq \frac{50\,\text{V}}{R_t} \]

dove:  \\
- Idn è la corrente differenziale nominale del dispositivo di protezione. \\
- 50 V è la massima tensione ammissibile sulle masse. \\

Rt è la resistenza del dispersore e dei conduttori di protezione delle masse.
L’impianto elettrico considerato nel progetto seguente, trattandosi di un impianto domestico, è distribuito mediante l’utilizzo del sistema TT. \\
Per l’alimentazione delle prese sono stati utilizzati interruttori differenziali con una corrente nominale differenziale di 30 mA, di tipo “A”, per correnti sinusoidali e pulsanti, in conformità alla norma 64-8, poiché all'interno della Wallbox è integrata una protezione contro le correnti continue di 6mA. \\
Considerando un interruttore differenziale con corrente nominale di intervento pari a 30 mA (Idn), in riferimento alla disequazione riportata in precedenza, la resistenza di terra deve risultare inferiore o uguale a 1666,67 $\Omega$. In caso contrario, l’impianto di terra dovrà necessariamente essere adeguato.

\subsection{Protezione dai contatti diretti}
La protezione contro i contatti diretti deve essere conforme a quanto previsto nella norma CEI 64–8 artt. 412.1, 412.2: saranno pertanto utilizzate “misure di protezione totale” mediante isolamento delle parti attive e/o mediante involucri o barriere. \\
\textbf{Si evidenzia che vernici, lacche, smalti e prodotti similari da soli non sono considerati idonei per assicurare un adeguato isolamento per la protezione dai contatti diretti.} \\
In particolare, dovranno essere rispettate le seguenti prescrizioni: 
\begin{itemize}
    \item Parti attive ricoperte completamente con isolamento che può essere rimosso solo a
mezzo di distruzione;
    \item Grado di protezione minimo degli involucri contenenti parti attive dovrà essere almeno
IPXXB. Le superfici orizzontali degli involucri ubicati a portata di mano dovranno avere
grado di protezione non inferiore a IPXXD;
    \item Barriere o involucri devono poter essere rimossi o aperti solo con l’uso di un attrezzo
speciale o di una chiave.
\end{itemize}


\subsection{Potere di interruzione delle apparecchiature}
Gli interruttori installati all’interno dei quadri elettrici, dovranno avere potere di interruzione superiore alla massima corrente di corto circuito presunta nel punto di installazione degli stessi, tenendo in considerazione la formula seguente: 
\[ I_{\mathrm{cc-max}} < P\mathrm{.d.}I \]
dove: \\
- Icc-max è la corrente di cortocircuito massima.  \\
- P.d.I. è il potere di interruzione del dispositivo di protezione. \\
- Il potere di interruzione preso come riferimento nel seguente intervento progettuale è il potere di interruzione estremo (Icu) secondo la norma CEI EN 60947-2. \\


\subsection{Quadri elettrici}
\subsubsection{Riferimenti normativi}
L'installazione, l'assemblaggio e la posa in opera dei quadri elettrici devono essere eseguiti a regola d'arte, in conformità al D.M. 37/08 e nel rigoroso rispetto delle seguenti norme tecniche nelle edizioni in vigore all'atto dell'esecuzione dei lavori:


\begin{itemize}
\item \textbf{CEI 64-8 (IX Edizione 2024)}: Impianti elettrici utilizzatori a tensione nominale non superiore a 1000 V c.a. e a 1500 V c.c. (in particolare Parte 5 "Scelta ed installazione dei componenti elettrici" e Parte 6 "Verifiche").
\item \textbf{CEI EN 61439-1 e 61439-2:} Apparecchiature assiemate di protezione e di manovra per bassa tensione (Quadri di potenza). I quadri dovranno essere conformi anche alle indicazioni delle guide applicative CEI 17-113 e CEI 17-114.
\item \textbf{CEI EN 61439-3:} Quadri di distribuzione finale destinati a essere utilizzati da persone comuni (DBO).
\item \textbf{CEI 23-51:} Prescrizioni per quadri di distribuzione per installazioni fisse per uso domestico e similare (ove applicabile per tipologia di impianto).
\item \textbf{CEI 11-27 (VI Edizione 2025):} Lavori su impianti elettrici (sicurezza durante le fasi di installazione, manutenzione e collaudo).
\item \textbf{Regolamento UE n. 305/2011 (CPR):} Prodotti da costruzione (requisiti di reazione al fuoco per i cavi).
\end{itemize}

\subsubsection{Requisiti Costruttivi e di Assemblaggio}
La realizzazione del quadro deve garantire ordine, accessibilità e stabilità meccanica, rispettando i seguenti requisiti:
\begin{itemize}
\item \textbf{Modalità di fissaggio delle apparecchiature:}
\begin{itemize}
\item \textbf{Apparecchiature modulari:} Tutti i dispositivi con attacco standard (es. magnetotermici, differenziali, relè) devono essere fissati su guide profilate DIN (conformi alla EN 60715), garantendo un bloccaggio sicuro.
\item \textbf{Apparecchiature non modulari:} I dispositivi di altra tipologia (es. interruttori scatolati, sezionatori di potenza) devono essere installati su pannelli di fondo o piastre fissate su idonee traverse di sostegno, dimensionate per reggere il peso e le sollecitazioni meccaniche durante le manovre.
\end{itemize}

\item \textbf{Specifiche dei conduttori di cablaggio (CPR):}

Il cablaggio interno di potenza e comando deve essere realizzato con conduttori unipolari senza guaina. È fatto divieto di utilizzare cavi non conformi alla normativa vigente. I cavi devono possedere le seguenti caratteristiche:
\begin{itemize}
\item \textbf{Tipologia:} Cavo tipo FS17 con tensione nominale 450/750\,V (o superiore se richiesto dalla tensione di esercizio).
\item \textbf{Classe di reazione al fuoco:} Minimo Cca-s3, d1, a3 in conformità al Regolamento CPR.
\end{itemize}

\item \textbf{Dimensionamento e posa dei cablaggi:}
\begin{itemize}
\item \textbf{Sollecitazioni:} Le sezioni dei cavi e delle barre devono essere dimensionate per sopportare le sollecitazioni termiche (calore generato dalla corrente nominale $I_n$ in servizio continuo) e dinamiche (sforzi elettrodinamici derivanti dalla massima corrente di cortocircuito $I_{cc}$ presunta nel punto di installazione).
\item \textbf{Ordine:} I conduttori devono essere ordinati, fascettati ove necessario e disposti in modo da non esercitare trazioni meccaniche sui morsetti.
\item \textbf{Serraggio:} Il serraggio dei morsetti deve rispettare le coppie (N\,m) indicate dai costruttori, utilizzando preferibilmente chiavi dinamometriche per evitare surriscaldamenti o allentamenti nel tempo.
\end{itemize}
\end{itemize}


\subsubsection{Prescrizioni generali di posa in opera}
L'installatore è tenuto a rispettare le seguenti prescrizioni operative per garantire la sicurezza, la manutenibilità e la durabilità dell'opera:

\begin{itemize}
    \item \textbf{Ubicazione e accessibilità}:
    \begin{itemize}
        \item \textit{Spazi di rispetto}: Il quadro deve essere installato in posizione accessibile. Davanti al quadro deve essere mantenuta una distanza libera da ostacoli non inferiore a \textbf{70 cm} (o sufficiente a permettere l'apertura completa delle portelle a 90° o 180°).
        \item \textit{Altezza di installazione}: I dispositivi di comando e sezionamento devono trovarsi ad un'altezza compresa preferibilmente tra \textbf{0,6 m e 1,7 m} dal piano di calpestio, per consentire un agevole intervento in caso di emergenza.
        \item \textit{Illuminazione}: L'area di installazione deve essere adeguatamente illuminata, anche tramite illuminazione di sicurezza ove previsto, per l'identificazione certa dei circuiti.

    \end{itemize}
    \item \textbf{Fissaggio Meccanico e Grado di Protezione (IP)}:
    \begin{itemize}
    \item \textbf{Ancoraggio:} Il fissaggio a parete deve essere solido, realizzato mediante tasselli adeguati alla tipologia di muratura e garantendo la perfetta verticalità dell'assieme.
    
    \item \textbf{Mantenimento del grado IP:} L'ingresso delle condutture deve avvenire esclusivamente tramite appositi bocchettoni, pressacavi o flange che garantiscano il mantenimento del grado di protezione (IP) di targa del quadro. È vietato lasciare fori aperti o non sigillati.
    
    \item \textbf{Integrità dell'involucro:} Dopo l'installazione, devono essere ripristinate tutte le barriere (pannelli frontali, mostrine) per garantire il grado di protezione contro i contatti diretti (minimo \textbf{IPXXB} o \textbf{IP2X} a portella aperta, \textbf{IP4X} per superfici orizzontali accessibili).
    \end{itemize}
\end{itemize}


\subsubsection{Identificazione, Documentazione e Verifiche Finali}
Al termine dell'installazione, l'installatore deve provvedere alle seguenti attività obbligatorie:

\begin{itemize}
\item \textbf{Identificazione:} Ogni interruttore e dispositivo di comando deve essere chiaramente identificato tramite apposita targhetta o etichetta indelebile, riportante la funzione o il circuito comandato, in coerenza con gli schemi unifilari. Sull'involucro deve essere applicata la targa identificativa del costruttore (conformità CEI EN 61439) e il segnale di pericolo ``Tensione Elettrica''.

\item \textbf{Documentazione:} Inserire nella tasca porta-documenti (se presente) o consegnare alla committenza lo schema unifilare aggiornato (``as-built'').

\item \textbf{Verifiche strumentali (CEI 64-8/6):} Prima della messa in servizio, devono essere eseguite le verifiche iniziali con strumentazione idonea:
\begin{enumerate}
\item \textit{Continuità del PE:} Verifica della continuità del conduttore di protezione verso la barra di terra del quadro (corrente di prova $\geq 200$\,mA).
\item \textit{Isolamento:} Misura della resistenza di isolamento (valore minimo $\geq 1$\,M$\Omega$ a 500\,Vcc).
\item \textit{Differenziali:} Prova di scatto degli interruttori differenziali (tempo e corrente d'intervento) e prova funzionale tasto ``Test''.
\end{enumerate}
\end{itemize}

L'esito di tali prove dovrà essere riportato nel verbale di verifica allegato alla Dichiarazione di Conformità o ne va dichiarata la conformità.


\subsubsection{Caratteristiche elettriche}

\begin{center}
\begin{tabular}{l@{\hspace{2em}}l}
Tensione nominale:   & 230\,V \\
Numero di fasi:      & F+N \\
Frequenza nominale: & 50\,Hz \\
Isolamento:          & in aria \\
Ventilazione:        & naturale \\
\end{tabular}
\end{center}


